\documentclass{handout}

\assignment{Homework 9}
\title{Ray optics}
\duedate{Wednesday November 19, 2025}

\begin{document}
\maketitle

\hwinstructions

\begin{questions}
    \question \textbf{Speed of light}: 	There was a major collision of an asteroid with the Moon in medieval times. It was described by monks at Canterbury Cathedral in England as a red glow on and around the Moon. How long after the asteroid hit the Moon, which is  $3.84\times10^5$ km away, would the light first arrive on Earth?

    \question \textbf{Reflection}: 	Suppose a man stands in front of a mirror as shown below. His eyes are 1.65 m above the floor and the top of his head is 0.13 m higher. Find the height above the floor of the top and bottom of the smallest mirror in which he can see both the top of his head and his feet. How is this distance related to the man’s height?
    \begin{figure}[h!]
        \centering
        \includegraphics[height=2in]{Images/mirror.png}
    \end{figure}

    \question \textbf{Scuba perspectives}: A scuba diver who is 2.0 m below the surface looks up and sees his diving instructor standing on the edge of the water. This ray (in the water) forms an angle of 25$^\circ$ with the surface of the water. Water has a refractive index of 1.33 (4/3).
    \begin{parts}
        \part What is the angle the ray forms with the vertical outside the water?
        \part The location where the ray passes the interface is 2 m from the edge of the water. How high is the instructor's head above the surface of the water?
        \part Explain why the diver perceives the instructor to be higher than they actually are.
        \part Explain why the instructor perceives the diver to be shallower than they actually are.
    \end{parts}

    \clearpage
    \question \textbf{Images in a plane mirror}: Consider a pair of flat mirrors that are positioned so that they form an angle of 60$^\circ$. An object is placed on the bisector between the mirrors. Construct an analogous ray diagram to below to determine how many images would be formed. A schematic of the mirror is given. Draw virtual rays coming from each image as well as the corresponding real rays originating from the object.
    \begin{figure}[h!]
        \centering
        \includegraphics[width=0.4\linewidth]{Images/twoMirrorExample.png}
        \hspace{.5in}
        \includegraphics[width=0.4\linewidth]{Images/mirrorproblem.pdf}
    \end{figure}

    \question \textbf{Convex lens}: You have a convex magnifying glass with a focal length of 6 in. Draw ray diagrams and use the thin lens equation.
    \begin{parts}
        \part If the object is 10 in from the glass, where is the image located, and is it virtual or real?
        \part If the object is 2 in away, where is the image located, and is it virtual or real?
    \end{parts}
    \begin{figure}[h!]
        \centering
        \includegraphics[width=0.7\linewidth]{Images/lenses.pdf}
    \end{figure}
\end{questions}

\end{document}
